\section{Euclidean Scale and Shape}

Section~3 established the orthogonal decomposition
\[
\mathbf h
=
\bar h\,\mathbf 1+\mathbf r,
\qquad
\sum_{i=1}^{6}r_i=0.
\]

This section proves its geometric meaning for the fixed-normal support
realization
\[
P(\mathbf h)
=
\left\{
\mathbf x\in\mathbb E^3:
|\mathbf n_i\cdot\mathbf x|\le h_i
\text{ for all }i
\right\}.
\]

The uniform direction is exactly the regular Euclidean scale direction.
Within the dodecahedral support chamber, a nonzero residual gives a
non-homothetic shape variation.

\subsection{Support scaling}

\begin{theorem}[Support scaling]
For every
\[
\mathbf h\in\mathbb R_{>0}^6
\]
and every \(c>0\),
\[
P(c\mathbf h)=cP(\mathbf h).
\]
\end{theorem}

\begin{proof}
Let \(\mathbf x\in P(c\mathbf h)\) and set
\[
\mathbf y=\frac{1}{c}\mathbf x.
\]
Then
\[
|\mathbf n_i\cdot\mathbf y|
=
\frac{1}{c}|\mathbf n_i\cdot\mathbf x|
\le h_i
\]
for every \(i\). Hence
\[
\mathbf y\in P(\mathbf h),
\]
so
\[
\mathbf x\in cP(\mathbf h).
\]

Conversely, if
\[
\mathbf x=c\mathbf y
\]
with \(\mathbf y\in P(\mathbf h)\), then
\[
|\mathbf n_i\cdot\mathbf x|
=
c|\mathbf n_i\cdot\mathbf y|
\le ch_i.
\]
Thus
\[
\mathbf x\in P(c\mathbf h).
\]
\end{proof}

This result holds throughout the positive support domain and does not
require preservation of dodecahedral incidence.

\subsection{The regular similarity family}

Let
\[
D=P(h_0\mathbf 1)
\]
be the reference regular dodecahedron.

For every \(h>0\),
\[
P(h\mathbf 1)
=
\frac{h}{h_0}D.
\]

Therefore the positive uniform ray
\[
U_+
=
\{h\mathbf 1:h>0\}
\]
realizes the regular origin-centered similarity family with the fixed
normal directions.

\subsection{Uniqueness of the regular scale direction}

For a compact convex set \(K\subset\mathbb E^3\), define its support value
in the unit direction \(\mathbf u\) by
\[
s_K(\mathbf u)
=
\max_{\mathbf x\in K}
\mathbf u\cdot\mathbf x.
\]

If the plane
\[
\mathbf n_i\cdot\mathbf x=h_i
\]
supports a genuine facet of \(P(\mathbf h)\), then
\[
s_{P(\mathbf h)}(\mathbf n_i)=h_i.
\]

\begin{theorem}[Uniqueness of the regular scale direction]
Let
\[
\mathbf h\in\mathbb R_{>0}^6,
\]
and assume that all twelve defining planes of \(P(\mathbf h)\) support
genuine facets. If
\[
P(\mathbf h)=cD
\]
for some \(c>0\), then
\[
\mathbf h=ch_0\mathbf 1.
\]
\end{theorem}

\begin{proof}
For every \(i\),
\[
h_i
=
s_{P(\mathbf h)}(\mathbf n_i).
\]

Since support values scale under homothety,
\[
s_{P(\mathbf h)}(\mathbf n_i)
=
c\,s_D(\mathbf n_i)
=
ch_0.
\]

Thus
\[
h_i=ch_0
\]
for all \(i\), and hence
\[
\mathbf h=ch_0\mathbf 1.
\]
\end{proof}

The active-facet assumption excludes redundant inequalities whose written
support values are not recoverable from the realized polytope.

\subsection{Residual variation}

Write
\[
\mathbf h
=
\bar h\,\mathbf 1+\mathbf r,
\qquad
\mathbf r\in H.
\]

If
\[
\mathbf r=\mathbf 0,
\]
then
\[
\mathbf h=\bar h\,\mathbf 1,
\]
and \(P(\mathbf h)\) is regular.

If
\[
\mathbf r\ne\mathbf 0,
\]
then the six support numbers are not all equal. Under the active-facet
condition, the uniqueness theorem shows that \(P(\mathbf h)\) cannot be a
homothetic copy of the regular dodecahedron.

\begin{corollary}[Residual criterion]
Assume that all twelve defining planes remain facet-defining. Then
\[
\mathbf r=\mathbf 0
\]
if and only if \(P(\mathbf h)\) lies in the regular similarity family.
\end{corollary}

\subsection{Local preservation of dodecahedral incidence}

Positive support numbers do not alone guarantee the regular dodecahedral
combinatorial type. The required stability holds locally.

\begin{theorem}[Local stability of dodecahedral incidence]
There exists an open neighborhood
\[
\mathcal U\subset\mathbb R_{>0}^6
\]
of \(h_0\mathbf 1\) such that every \(P(\mathbf h)\) with
\[
\mathbf h\in\mathcal U
\]
has the same face-vertex incidence structure as the regular dodecahedron.
\end{theorem}

\begin{proof}
Write the twelve signed unit normals as
\[
\mathbf m_1,\ldots,\mathbf m_{12},
\]
and let \(p(a)\in\{1,\ldots,6\}\) identify the opposite-face pair
containing \(\mathbf m_a\).

Then
\[
P(\mathbf h)
=
\left\{
\mathbf x:
\mathbf m_a\cdot\mathbf x
\le
h_{p(a)}
\text{ for }a=1,\ldots,12
\right\}.
\]

The regular dodecahedron is simple, so each vertex is determined by an
invertible triple
\[
T=\{a,b,c\}
\]
of incident facet planes. Let
\[
N_T
=
\begin{pmatrix}
\mathbf m_a^{\mathsf T}\\
\mathbf m_b^{\mathsf T}\\
\mathbf m_c^{\mathsf T}
\end{pmatrix}.
\]

The corresponding intersection point is
\[
\mathbf v_T(\mathbf h)
=
N_T^{-1}
\begin{pmatrix}
h_{p(a)}\\
h_{p(b)}\\
h_{p(c)}
\end{pmatrix},
\]
which depends continuously on \(\mathbf h\).

At the regular state, every nonincident facet inequality is strict at each
regular vertex. Because only finitely many such inequalities occur, their
minimum slack is positive. Hence all regular vertex triples remain feasible
throughout a sufficiently small neighborhood.

Conversely, every invertible triple that is not a regular vertex violates
at least one defining inequality strictly at the regular state. Finiteness
and continuity allow the neighborhood to be chosen so that no such triple
becomes feasible.

After also preserving positivity, exactly the original vertex triples
remain feasible. Therefore the vertex-facet incidence structure, and hence
the full dodecahedral combinatorial type, is unchanged.
\end{proof}

\subsection{The dodecahedral support chamber}

\begin{definition}[Dodecahedral support chamber]
Let \(\mathcal C_D\) be the connected component containing
\(h_0\mathbf 1\) of the set of positive support registers for which

\begin{enumerate}
\item all twelve defining planes support genuine facets;
\item the polytope is simple; and
\item its face-vertex incidence structure is dodecahedral.
\end{enumerate}
\]
\end{definition}

The local stability theorem guarantees that \(\mathcal C_D\) contains an
open neighborhood of the regular state.

No claim is made that
\[
\mathcal C_D=\mathbb R_{>0}^6.
\]

Its boundary may contain facet loss, edge contraction, vertex collision, or
other changes of combinatorial type.

\subsection{Euclidean scale-shape theorem}

\begin{theorem}[Euclidean scale-shape decomposition]
Let
\[
\mathbf h\in\mathcal C_D
\]
and write
\[
\mathbf h
=
\bar h\,\mathbf 1+\mathbf r,
\qquad
\sum_{i=1}^{6}r_i=0.
\]

Then:

\begin{enumerate}
\item \(\bar h\,\mathbf 1\) is the unique regular support register with mean
support \(\bar h\);

\item for every \(c>0\),
\[
P(c\mathbf h)=cP(\mathbf h);
\]

\item \(P(\mathbf h)\) is regular if and only if
\[
\mathbf r=\mathbf 0;
\]

\item if
\[
\mathbf r\ne\mathbf 0,
\]
then \(P(\mathbf h)\) is not an origin-centered homothetic copy of \(D\).
\end{enumerate}
\end{theorem}

\begin{proof}
The first statement follows from the regular uniform family and uniqueness
of the regular scale direction. The second is the support-scaling theorem.

If \(\mathbf r=\mathbf 0\), then the support register is uniform and the
realization is regular. Conversely, a regular realization in
\(\mathcal C_D\) must have equal support values, so its residual vanishes.
The final statement follows immediately.
\end{proof}

Thus
\[
\text{uniform component}
\longleftrightarrow
\text{regular Euclidean scale},
\]
while
\[
\text{nonzero zero-mean residual}
\longleftrightarrow
\text{non-homothetic support shape variation}.
\]

\subsection{Boundary}

The theorem applies only to the fixed-normal, origin-centered, centrally
symmetric realization family and only inside \(\mathcal C_D\).

It does not classify arbitrary Euclidean dodecahedra or arbitrary
deformations of the regular solid. It excludes varying normal directions,
unequal opposite supports, translations, independent face rotations, and
arbitrary vertex motion.

The terms \emph{scale} and \emph{shape} are used only in their ordinary
Euclidean sense. No mechanics, elasticity, force, energy, dynamics, or
physical expansion is implied.
